% =====================================================================
%  Literature Review (standalone) --- prepared for supervisory review
%  Project: SE(3)-equivariant recurrence for viewpoint-invariant,
%           edge-deployable rehabilitation assessment
%  Compile:  pdflatex literature_review_egru && pdflatex literature_review_egru
%  NOTE: standalone document, NOT the paper's condensed Related Work section.
% =====================================================================
\documentclass[11pt]{article}
\usepackage[margin=1in]{geometry}
\usepackage{amsmath,amssymb}
\usepackage[hidelinks]{hyperref}
\usepackage{cleveref}
\setlength{\parskip}{0.35em}

\title{\vspace{-2em}Literature Review: Geometric Guarantees, Node Robustness, and\\
Protocol Rigor in Skeleton-Based Rehabilitation Assessment}
\author{}
\date{}

\begin{document}
\maketitle
\vspace{-3em}

\begin{abstract}
\noindent Automated scoring of rehabilitation exercises from skeleton data is maturing as a
benchmark task, but the conditions under which it is evaluated diverge sharply from the conditions
under which a home monitoring device operates. This review surveys three bodies of work that bear on
that gap: (i) clinical movement-analysis models and the evaluation protocols that rank them; (ii)
equivariant architectures and their behaviour under sensor-node failure; and (iii) frame
canonicalization as an alternative route to viewpoint invariance, contrasted with architectural
guarantees under causal streaming. In each case we characterise what prior work establishes, where it
leaves a measurable gap, and how a certified-invariant, edge-deployable formulation addresses it. The
recurring theme is a distinction between properties that hold \emph{in aggregate on a fixed corpus} and
properties that are \emph{guaranteed per prediction and off distribution} --- a distinction that is
decisive for a clinical instrument and largely absent from the current literature.
\end{abstract}

% =====================================================================
\section{Clinical Movement Analysis and Protocol Rigor}
Skeleton-based assessment of physical rehabilitation has consolidated around a small number of public
corpora, chief among them KIMORE~\cite{capecci2019kimore}, which pairs Kinect-tracked exercise
recordings with clinician-assigned scores and has become the de facto ranking ground for the task.
Progress on it is reported almost exclusively through supervised backbones of increasing capacity:
spatial--temporal graph convolutional networks (ST-GCN)~\cite{yan2018stgcn} and their clinical
descendants trained with supervised-contrastive objectives~\cite{karlov2024}, cross-modal video-to-joint
augmentation to enlarge scarce training data~\cite{abedi2023}, dual-stream graph networks with
motion-aware grouping~\cite{kuang2026}, and, most recently, point-cloud transformers that treat the
skeleton as an unordered set and attend over it~\cite{pct2026}. These models differ in inductive bias
and parameter count, but they share an evaluation contract: train and test on clean, frontal recordings
from a single sensor, and rank by a single aggregate error against the clinician's score.

Two properties of this contract limit what its leaderboard can be taken to mean, and both are
under-examined in the literature. \textbf{First}, none of these architectures carries an explicit
geometric guarantee under a change of camera. Viewpoint robustness, where it is claimed at all, is an
emergent, aggregate property obtained by training-time rotation augmentation; it says nothing about
whether the score assigned to a \emph{particular} patient is stable when the same performance is filmed
from a different wall. Aggregate invariance and per-prediction invariance are not the same property, and
the field's metrics cannot distinguish them. \textbf{Second}, the reporting protocol on the most
commonly used single-exercise slice is vulnerable to optimistic selection: when the reported epoch is
chosen by its minimum on the test set, and the input scaler is fit over the full dataset, the headline
number is inflated relative to an honest, validation-selected estimate. Rehab-Pile~\cite{rehabpile2026},
which aggregates KIMORE, UI-PRMD~\cite{vakanski2018uiprmd} and IRDS into a cross-subject benchmark, is
the closest antecedent of a more disciplined protocol, but even it does not report the one reference a
degradation curve requires for interpretation: a mean-predictor floor. Absent that floor, a model that
ignores its input entirely --- and therefore has zero degradation under any corruption --- would win
every robustness plot ever drawn.

The gap, then, is methodological before it is architectural. A rigorous treatment must (a) establish a
noise floor below which no accuracy difference is interpretable, since training on this task is not
bit-reproducible and repeated fits of an identical configuration vary by a non-trivial margin; (b)
quantify protocol inflation directly rather than assume it away; and (c) attach uncertainty to
``beats the floor'' claims with a resampling scheme whose unit is the \emph{subject}, not the sequence,
so that within-subject correlation is not mistaken for signal. Our work adopts this stance: we measure
the run-to-run nondeterminism floor ($\approx\!0.33$~MAD), show that optimistic epoch selection is worth
an $18.6\%$ MAD reduction on the single-exercise slice, and demonstrate with a paired subject-cluster
bootstrap that the honest number on that slice does not reliably separate from a mean predictor. Every
subsequent claim is pooled, subject-disjoint, and gated on clearing the floor --- so that benchmark
noise is never reported as clinical signal.

% =====================================================================
\section{Equivariant Architectures and Node Robustness}
The alternative to inducing invariance through augmentation is to build it into the hypothesis class.
Group-equivariant networks constrain a model to commute with a symmetry group by
construction~\cite{cohen2016gconv}, so that the guarantee holds for arbitrary weights rather than being
learned. For the Euclidean group in three dimensions, two design philosophies dominate. The
\emph{steerable} family --- tensor-field networks~\cite{thomas2018tfn}, $SE(3)$-Transformers with
equivariant attention~\cite{fuchs2020se3transformer}, and the \texttt{e3nn} framework that systematises
them~\cite{geiger2022e3nn} --- represents features in irreducible representations of $O(3)$ and couples
them through spherical-harmonic tensor products, yielding exact equivariance at the cost of additional
machinery. The \emph{lightweight} family instead reaches $E(n)$-equivariance without spherical harmonics:
EGNN~\cite{satorras2021egnn} updates node states from relative displacements $x_i-x_j$ and their norms,
and vector-neuron constructions~\cite{deng2021vectorneurons} lift standard operators to $SO(3)$-equivariant
ones on vector-valued features. The lightweight route is attractive precisely because it dispenses with
the steerable overhead while retaining a rotation guarantee.

The gap this review highlights concerns \emph{robustness to sensor failure}, which is not addressed by
the equivariance property itself and separates these families in practice. Relative-coordinate updates
couple every node's representation to its neighbours' coordinates; when a joint is corrupted --- a frozen
or occluded Kinect node, the characteristic failure mode of consumer depth cameras --- the corruption
enters the message graph and propagates. This is orthogonal to whether the network is equivariant: a
model can be exactly rotation-invariant and still fragile to a single dead joint. The prevailing
responses in the skeleton literature --- masking or attention over occluded joints~\cite{yan2018stgcn}
--- treat failure as noise to be suppressed rather than as a discriminating test between architectures
that are otherwise indistinguishable on clean data.

Our position is that node robustness, not viewpoint invariance, is what justifies the heavier steerable
encoder, and that the justification must be \emph{measured} rather than asserted. We are explicit that
steerable machinery is not necessary for viewpoint invariance --- a recurrence over six hand-crafted
$SO(3)$-invariant features attains the same accuracy and the same guarantee at a fraction of the
parameters. What separates the models is a single axis: behind an \emph{identical} invariant read-out,
the lighter $E(n)$ update is empirically the more brittle of the two, crossing the mean-predictor floor
at two dead joints where the steerable, graph-pooled encoder still degrades gracefully; and the
named-feature baseline, whose features explicitly name the joints they depend on, collapses through the
floor at a \emph{single} dead joint. The steerable representation earns its place not by a stronger
theorem but by diluting a local failure across the learned message set --- a claim we support with
controlled dead-joint experiments rather than a robustness proof, in keeping with the distinction between
what is guaranteed and what is merely observed.

% =====================================================================
\section{Frame Canonicalization versus Architectural Guarantees in Causal Streaming}
A third route to viewpoint invariance sidesteps equivariant architectures entirely: align each input to
a canonical reference frame, then apply an ordinary non-equivariant backbone. Principal-axis (PCA)
alignment is the classical instance, and frame averaging~\cite{puny2022frameaveraging} generalises it by
symmetrising a network's output over a set of frames rather than committing to one. These heuristics are
genuinely competitive: on our benchmark a PCA-canonicalized transformer matches the equivariant model on
every clean, offline axis, and we report this concession directly rather than obscure it. Any honest
comparison must begin by conceding that canonicalization is a strong baseline, not a straw man.

The gap is in the \emph{kind} of invariance a canonical frame provides. A PCA frame is a point estimate
of an orientation, obtained from the spatial covariance of the joints; it is exact only where that
covariance is well-conditioned. On clinical poses it frequently is not --- limb extensions and other
planar configurations drive the pose covariance toward degeneracy, at which point the eigenvector
ordering and sign convention that define the frame become discontinuous. We measure this on KIMORE:
$18\%$ of frames sit within an eigen-gap of $0.05$ of degeneracy, and the estimated frame rotates by more
than $45^\circ$ across $21\%$ of frame transitions, up to full sign flips. A downstream network can learn
to absorb this jitter \emph{on a given corpus}, but no theorem bounds the behaviour off distribution ---
which is exactly the regime a deployed instrument enters when it meets a camera placement it never saw in
training. We note, against a common misconception, that the per-frame spatial PCA frame is itself
\emph{causal} and requires no temporal lookahead; the difficulty is not latency but the absence of a
guarantee and the presence of covariance-driven discontinuities.

The distinction our work draws is therefore between a heuristic that is strong on the corpus it was tuned
on and an architectural guarantee that transfers. Our read-out estimates no frame at all: it is an
explicit projection onto a generating set of $O(3)$ invariants, exact for arbitrary weights and certified
numerically to the fp64 roundoff floor ($\sim\!10^{-16}$). Because the guarantee is a property of the
architecture rather than of the data, it transfers zero-shot to a genuinely novel camera --- the real
$45^\circ$ NTU~RGB+D Cross-View displacement --- without retraining and without a per-frame covariance
estimate that can degenerate. It also composes cleanly with the deployment path: the read-out converts to
a strictly causal streaming form that emits a first score roughly $8.9$~ms after the first frame, and the
guarantee survives int8 quantization, so viewpoint invariance is retained on the edge rather than
purchased offline. Where the field has treated irregular sampling and continuous time as an additional
frontier~\cite{chen2018neuralode,rubanova2019odernn,kidger2020neuralcde,che2018grud}, we find --- and
report as a negative result, quantified with a spectral analysis of the signal
band~\cite{vanderplas2018lombscargle} and implemented through an adaptive Dormand--Prince
solver~\cite{dormand1980prince} --- that native continuous-time consumption buys nothing on this corpus,
reinforcing that the payoff lies in the geometric guarantee, not in the temporal machinery.

% =====================================================================
\section{Synthesis: The Gap and Our Position}
Read together, these three literatures leave a coherent gap. Clinical assessment models compete on
aggregate error under a protocol that can inflate the headline number and reports no floor against which
degradation is interpretable. Equivariant architectures supply a rotation guarantee but say nothing about
sensor failure, and the lightweight variants that are cheapest to deploy are also the most fragile to it.
Canonicalization offers strong offline invariance but only as an uncertified point estimate that turns
discontinuous on exactly the degenerate poses clinical exercises produce. In every case the missing
ingredient is the same: a property that is \emph{guaranteed per prediction and holds off the training
distribution}, rather than one that holds in aggregate on a fixed corpus.

The work this review supports occupies that gap directly. It treats viewpoint invariance as a certified
theorem of the architecture rather than an emergent effect of augmentation; it justifies the steerable
encoder on the one measured axis --- sensor-node robustness --- where it demonstrably differs from lighter
alternatives; it confronts canonicalization on its own terms, conceding the offline tie and locating the
difference in the guarantee; and it carries all of this onto the edge under a protocol disciplined enough
that benchmark noise cannot be mistaken for clinical signal. The contribution is less a new backbone than
a reframing of what the benchmark should be asked to measure, backed by guarantees and controls the
current literature does not provide.

% =====================================================================
\begin{thebibliography}{00}
\bibitem{capecci2019kimore} M.~Capecci et al., ``The KIMORE dataset: Kinematic assessment of movement and clinical scores for remote monitoring of physical rehabilitation,'' \textit{IEEE TNSRE}, vol.~27, no.~7, pp.~1436--1448, 2019.
\bibitem{pct2026} ``A point cloud transformer for remote monitoring and automated assessment of physical rehabilitation exercises,'' arXiv:2606.30309, 2026.
\bibitem{karlov2024} M.~Karlov, A.~Abedi, and S.~S.~Khan, ``Rehabilitation exercise quality assessment through supervised contrastive learning with hard and soft negatives,'' \textit{Med. Biol. Eng. Comput.}, 2024.
\bibitem{abedi2023} A.~Abedi, M.~Malmirian, and S.~S.~Khan, ``Cross-modal video to body-joints augmentation for rehabilitation exercise quality assessment,'' arXiv:2306.09546, 2023.
\bibitem{kuang2026} Z.~Kuang et al., ``Dual-stream STGCN with motion-aware grouping for rehabilitation action quality assessment,'' \textit{Sensors}, vol.~26, no.~1, p.~287, 2026.
\bibitem{rehabpile2026} A.~Ismail-Fawaz et al., ``Rehab-Pile: A cross-subject rehabilitation benchmark aggregating KIMORE, UI-PRMD and IRDS,'' arXiv:2507.21018, IEEE FG 2026.
\bibitem{vakanski2018uiprmd} A.~Vakanski et al., ``A data set of human body movements for physical rehabilitation exercises (UI-PRMD),'' \textit{Data}, vol.~3, no.~1, p.~2, 2018.
\bibitem{yan2018stgcn} S.~Yan, Y.~Xiong, and D.~Lin, ``Spatial temporal graph convolutional networks for skeleton-based action recognition,'' in \textit{Proc. AAAI}, 2018.
\bibitem{cohen2016gconv} T.~S.~Cohen and M.~Welling, ``Group equivariant convolutional networks,'' in \textit{Proc. ICML}, 2016.
\bibitem{thomas2018tfn} N.~Thomas et al., ``Tensor field networks: Rotation- and translation-equivariant neural networks for 3D point clouds,'' arXiv:1802.08219, 2018.
\bibitem{fuchs2020se3transformer} F.~B.~Fuchs et al., ``SE(3)-Transformers: 3D roto-translation equivariant attention networks,'' in \textit{Proc. NeurIPS}, 2020.
\bibitem{geiger2022e3nn} M.~Geiger and T.~Smidt, ``e3nn: Euclidean neural networks,'' arXiv:2207.09453, 2022.
\bibitem{satorras2021egnn} V.~G.~Satorras, E.~Hoogeboom, and M.~Welling, ``E(n) equivariant graph neural networks,'' in \textit{Proc. ICML}, 2021.
\bibitem{deng2021vectorneurons} C.~Deng et al., ``Vector neurons: A general framework for SO(3)-equivariant networks,'' in \textit{Proc. ICCV}, 2021.
\bibitem{puny2022frameaveraging} O.~Puny et al., ``Frame averaging for invariant and equivariant network design,'' in \textit{Proc. ICLR}, 2022.
\bibitem{chen2018neuralode} R.~T.~Q.~Chen et al., ``Neural ordinary differential equations,'' in \textit{Proc. NeurIPS}, 2018.
\bibitem{rubanova2019odernn} Y.~Rubanova, R.~T.~Q.~Chen, and D.~Duvenaud, ``Latent ODEs for irregularly-sampled time series,'' in \textit{Proc. NeurIPS}, 2019.
\bibitem{kidger2020neuralcde} P.~Kidger et al., ``Neural controlled differential equations for irregular time series,'' in \textit{Proc. NeurIPS}, 2020.
\bibitem{che2018grud} Z.~Che et al., ``Recurrent neural networks for multivariate time series with missing values,'' \textit{Sci. Rep.}, vol.~8, p.~6085, 2018.
\bibitem{dormand1980prince} J.~R.~Dormand and P.~J.~Prince, ``A family of embedded Runge--Kutta formulae,'' \textit{J. Comput. Appl. Math.}, vol.~6, no.~1, pp.~19--26, 1980.
\bibitem{vanderplas2018lombscargle} J.~T.~VanderPlas, ``Understanding the Lomb--Scargle periodogram,'' \textit{ApJS}, vol.~236, no.~1, p.~16, 2018.
\end{thebibliography}

\end{document}
